一般超胞由一个整数矩阵 $S$ 定义。矩阵的每一行说明一条新晶格矢量由原来的三条矢量怎样组合。例如，第一行为 $(2,1,0)$，就表示 $\mathbf a'=2\mathbf a+\mathbf b$。在本书的行向量约定下，晶格、体积和位点数满足
\begin{equation}
 A'=SA,\qquad V'=|\det S|V,\qquad N'=|\det S|N.
 \label{eq:supercell-count-volume}
\end{equation}
$S$ 中的数必须是整数，且行列式不能为零。$|\det S|$ 是新旧晶胞的体积比，也是独立位点数的倍数；$N$ 始终数模型位点，不数边界上额外绘出的原子映像。对于三个重复数，$S=\operatorname{diag}(n_a,n_b,n_c)$，式\eqref{eq:supercell-count-volume}就回到三个重复数相乘的规则。上面的实机练习取 $S=2I$，所以 $|\det S|=8$，8个位点变为64个。

下面用一个非对角矩阵作代数检查：
\[
 S=\begin{bmatrix}2&1&0\\0&1&0\\0&0&1\end{bmatrix},\qquad
 \begin{cases}
 \mathbf a'=2\mathbf a+\mathbf b,\\
 \mathbf b'=\mathbf b,\quad \mathbf c'=\mathbf c.
 \end{cases}
\]
其行列式为2。若原模型含8个位点，新模型应含16个位点，体积是原来的两倍；新外框可以倾斜，周期材料本身的局部排列仍来自原结构的复制。这与通过应变拉伸原子间距有区别：超胞选择更大的重复单元，不把每条原子间连接都拉长两倍。

在\ui{超胞倍率或变换矩阵}中，可以填写三个重复整数，也可以填写完整的 $3\times3$ 矩阵。使用一般矩阵前先写出三条新基矢，再检查行列式、预览外框和应用后的位点数，便不容易把“重新定向”和“增加了多少原子”混为一谈。输入\literalinput{1 1 1}时，矩阵为 $I$，模型不扩大；本工具仍会把等价坐标归入单位晶胞，适合将落在胞外的代表位置整理到同一胞中查看。
